% Ad Familiares 13.65 --- headnote
% ad-familiares-13-65, undated (placeholder 54 BC), Cilicia / Bithynia province

\textit{Cicero to Publius Silius, propraetor of Bithynia
and Pontus, undated --- placed by Perseus's tradition
shortly after \textit{Fam.} 13.64, in the season of
Cicero's proconsulship of Cilicia (51--50 BC). The man
recommended is P. Terentius Hispo, an equestrian
\textit{pro magistro}, the deputy manager of one of the
great tax-farming partnerships --- specifically the
\textit{societas scripturae}, which held the publican
contract on the pasturage-tax of the Asian provinces.
Hispo's task was to negotiate \textit{pactiones},
fixed-sum settlements, with each Greek city; his
professional standing depended on closing them.}

\textit{Cicero himself had failed to obtain a
\textit{pactio} from Ephesus and, characteristically of
this cluster of letters, leans on Silius to succeed where
he could not. The double commendation --- Hispo himself,
and the wider \textit{societas} behind him --- shows
Cicero's habitual care to reinforce private bonds with
multiple anchorages: the partnership is in his
\textit{fides}, most of the partners are personal
friends, the deputy manager is an intimate of long
standing. The closing assurance that nothing in Silius's
whole command would be more welcome to Cicero is the
recurring sign-off of the genre.}
